\section{One stochastic rate: time discretization and algorithm}\label{app: oneRate}
As anticipated before, at each iteration of the algorithm and for each particle, we need to simulate 3 processes
$f_t$, $r_t$, $V_t$, and compute 3 integrals $U_t$, $\Theta_t$, $\Lambda_t$.

\subsection{Processes: $f_t$ and $r_t$}
Let $\{t_j\}$ be a discretization of $[0,T]$ and consider a single particle.\\
In order to simulate $(f_{t_{j+1}}, r_{t_{j+1}} - r^0_{t_{j+1}})$ from $(f_{t_{j}}, r_{t_{j}} - r^0_{t_{j}})$, one can observe that, from (\ref{eq: hullWhiteSolution}):

\begin{align*}
r_{t_{j+1}}& - r^0_{t_{j+1}}  =  e^{-\kappa(t_{j+1}-t_j) }(r_{t_{j}} - r^0_{t_{j}})+\sigma_r\int_{t_j}^{t_{j+1}}e^{-\kappa(t_{j+1}-t) }dB_t^T \\
& - \frac{\sigma^2_r}{\kappa^2}e^{-\kappa t_{j+1} }\big( \sinh(\kappa t_{j+1})(1-e^{-\kappa(T-t_{j+1}) })-\sinh(\kappa t_{j})(1-e^{-\kappa(T-t_{j}) }) \big) \numberthis \label{eq: discretization_interestRate}
\end{align*}

and use the following approximation scheme for $f_t$:

\begin{align*}
f_{t_{j+1}}& = f_{t_{j}} \exp\bigg(  \sigma(t_j,P_{t_j T}f_{t_j})(W_{t_{j+1}}^T-W^T_{t_{j}})+\sigma_r\int_{t_j}^{t_{j+1}} \frac{1-e^{-\kappa(T-t)}}{\kappa}dB_t^T \\
& -\frac{1}{2}(t_{j+1}-t_{j})\big(\sigma(t_j,P_{t_j T}f_{t_j})^2+2\rho\sigma_r\sigma(t_j,P_{t_j T}f_{t_j})I^1_j + \sigma_r^2I_j^2\big)  \    \bigg) \numberthis \label{eq: discretization_forwardPrice}
\end{align*}
with 
\begin{align*}
& I^1_j = \frac{\kappa(t_{j+1}-t_{j})-e^{-\kappa T }(e^{\kappa t_{j+1} }-e^{\kappa t_{j} }) }{\kappa^2 (t_{j+1}-t_{j})} \\
& I^2_j = \frac{\kappa(t_{j+1}-t_{j})  -2e^{-\kappa T }(e^{\kappa t_{j+1} }-e^{\kappa t_{j} }) +\frac{1}{2}e^{-2\kappa T }(e^{2\kappa t_{j+1} }-e^{2\kappa t_{j} }) }{\kappa^3(t_{j+1}-t_{j})}
\end{align*}
Then the simulation of $(f_{t_{j+1}}, r_{t_{j+1}} - r^0_{t_{j+1}})$ boils down to simulate exactly the Gaussian vector
\[\bigg( W_{t_{j+1}}^T-W^T_{t_{j}},  \int_{t_j}^{t_{j+1}}e^{-\kappa(t_{j+1}-t) }dB_t^T  ,\int_{t_j}^{t_{j+1}} \frac{1-e^{-\kappa(T-t)}}{\kappa}dB_t^T \bigg)  \]
whose covariance matrix is $(t_{j+1}-t_{j})\Sigma_j$ where
\[\Sigma_j= \left( \begin{array}{ccc}
1 & \rho J_j^1 & \rho I_j^1 \\ 
\rho J_j^1 & J_j^2 & L_j \\ 
\rho I_j^1 & L_j & I^2_j \end{array} \right)\] 
with
\begin{align*}
& J_j^1 = \frac{1-e^{-\kappa(t_{j+1}-t_j) }}{\kappa(t_{j+1}-t_j)} \\
& J_j^2 = \frac{1-e^{-2\kappa(t_{j+1}-t_j) } }{2\kappa(t_{j+1}-t_j)} \\
& L_j = \frac{J_j^1}{\kappa}-e^{-\kappa(T + t_{j+1}) } \frac{ e^{2\kappa t_{j+1} }-e^{2\kappa t_{j} }  }{2\kappa^2(t_{j+1}-t_{j})}
\end{align*}
\subsection{Process: $V_t$}
For $V_t$ we shall start with a plain Euler discretization scheme. Refinements might follow.\\
The computation of the partial derivative $\partial_S\sigma$ in $V_t$ will be addressed in dependence on the chosen interpolation technique.
\subsection{Integrals: $U_t$, $\Theta_t$, $\Lambda_t$}
First note that $V_t$ is computed over the time grid $\{t_j\}$, hence it is a discrete vector of values.
Therefore we set $ V_t = (V_{t_{j+1}}+V_{t_{j}})/2 $ for all $t\in [t_j,t_{j+1}]$. Same reasoning applies for $\Theta_t$.\\
From (\ref{eq: malliavinWeights}) we have: 

\begin{align*}
\Theta_{t_{j+1}} & =  \int_0^{t_{j+1}}\frac{e^{-\kappa s}}{V_s}ds = \Theta_{t_{j}} + \int_{t_j}^{t_{j+1}}\frac{e^{-\kappa s}}{V_s}ds = 
\Theta_{t_{j}} + \frac{2}{(V_{t_{j+1}}+V_{t_{j}})}\int_{t_j}^{t_{j+1}}e^{-\kappa s}ds = \\
 & = \Theta_{t_{j}} - \frac{2}{(V_{t_{j+1}}+V_{t_{j}})}\frac{1}{\kappa}(e^{-\kappa t_{j+1}}-e^{-\kappa t_{j}} )\\ \\
\Lambda_{t_{j+1}} &= \int_0^{t_{j+1}}\sigma_re^{2\kappa s}\Theta_s ds = \Lambda_{t_{j}} + \int_{t_j}^{t_{j+1}}\sigma_re^{2\kappa s}\Theta_s ds = \Lambda_{t_{j}} + \frac{(\Theta_{t_{j+1}} +\Theta_{t_{j}} )}{2}\sigma_r\int_{t_j}^{t_{j+1}}e^{2\kappa s} =\\
& = \Lambda_{t_{j}} + \frac{(\Theta_{t_{j+1}} +\Theta_{t_{j}} )}{4\kappa}\sigma_r(e^{2\kappa t_{j+1}} - e^{2\kappa t_{j}})
\end{align*}
\\
A more careful analysis is needed for $U_t$ as it involves the integration of the local volatility function.
From (\ref{eq: malliavinWeights}) we have:
$$  U_{t_{j+1}} = \int_0^{t_{j+1}}e^{\kappa s}\frac{\sigma(s,S_s)}{V_s}ds =  U_{t_{j}} +  \int_{t_{j}}^{t_{j+1}}e^{\kappa s}\frac{\sigma(s,S_s)}{V_s}ds$$
From step 4 in the algorithm at section \ref{subsec: algorithm}, we have that $\sigma(t,K)\equiv \sigma_N(t_j,K)$ for all $t\in [t_j,t_{j+1}]$, hence we can rewrite the integral as: 
$$ U_{t_{j+1}} = U_{t_{j}} +  \frac{2}{(V_{t_{j+1}}+V_{t_{j}})}\int_{t_{j}}^{t_{j+1}}e^{\kappa s}\sigma(t_j,S_s)ds $$

\IfFileExists{volatility.jpg}{\centerline{\includegraphics[scale=0.5]{volatility.jpg}}}{\centerline{\fbox{\parbox{0.6\textwidth}{\centering\small [figure \texttt{volatility.jpg}: place the file next to this source]}}}}
\bigskip
Now we have two possibilities:
\begin{itemize}
\item approximate $\sigma(t_j,S_s)$ as $[ \sigma(t_j,S_{t_{j+1}}) + \sigma(t_j,S_{t_j}) ]/2$ or $\sigma(t_j,\frac{S_{t_{j+1}} +S_{t_{j}}}{2}) $ and analytically solve the remaining integral
\item apply numerical integration techniques to the integral 
\end{itemize}
The advantage of the first approach is that it is faster to compute and does not depend on the interpolation of the local volatility but it could imply non-negligible errors in the case $S_{t_j}$ is far from $S_{t_{j+1}}$.
The numerical integration might be slower to compute and one should pay attention at the interpolation algorithm used for the local volatility. However it provides a more accurate approximation of the integral, hence we may choose it above the former.

\subsubsection*{Trapezoidal rule for $U_t$}
If using a trapezoidal rule, $U_{t_{j+1}}$ reads:
$$ U_{t_{j+1}} = U_{t_{j}} +  \frac{(t_{j+1}-t_j)}{(V_{t_{j+1}}+V_{t_{j}})}\big[\sigma(t_j,S_{t_j})e^{\kappa t_j}+ \sigma(t_j,S_{t_{j+1}})e^{\kappa t_{j+1}} \big] $$

\subsubsection*{Cavalieri-Simpson rule for $U_t$}
If using a Cavalieri-Simpson rule, $U_{t_{j+1}}$ reads:
$$ U_{t_{j+1}} = U_{t_{j}} +  \frac{1}{3}\frac{(t_{j+1}-t_j)}{(V_{t_{j+1}}+V_{t_{j}})}\big[\sigma(t_j,S_{t_j})e^{\kappa t_j}+  4\sigma\big(t_j,\frac{S_{t_{j+1}} +S_{t_{j}}}{2}\big)e^{\kappa \frac{t_{j+1}+t_{j}}{2}}   +\sigma(t_j,S_{t_{j+1}})e^{\kappa t_{j+1}} \big] $$
\smallskip

\subsection{The Algorithm}\label{subsec: algorithm}
Let $\{t_j\}$ be a discretization of $[0,T]$:

\begin{itemize}
\item[1.] Initialize $j=0$ and set $\sigma_N(t,K)=\sigma_{Dup}(0,K)$ for all $t\in [t_0,t_1]$
\item[2.] Simulate the $N$ particles  $\{(f_t^i, r_t^i- r^0_t) \}_{1\leq i\leq N}$ from  $t_{j}$ to  $t_{j+1}$ using  (\ref{eq: discretization_interestRate}) and (\ref{eq: discretization_forwardPrice}) 
\item[3.] For each particle, simulate $V_{t_{j+1}}^i$ and compute the integrals $U_{t_{j+1}}^i$, $\Theta_{t_{j+1}}^i$, $\Lambda_{t_{j+1}}^i$ according to (\ref{eq: malliavinWeights})
\item[4.] Sort the particles according to the spot value. For $K\in G_{K,t_{j+1}}$, find the smallest index $i(K)$ and the largest $j(K)$ for which $\delta_{t,N}(S_t^i-K) > \eta$. Compute the local volatility $\sigma_N(t_{j+1},K)$ from (\ref{eq: discrete_localVolatility_forward_HullWhite}).
\begin{itemize}
\item Interpolate it and extrapolate it
\item Set $\sigma(t,K)\equiv \sigma_N(t_{j+1},K)$ for all $t\in [t_{j+1},t_{j+2}]$
\end{itemize}
\item[5.] Set $j:=j+1$, iterate steps 2,3 and 4 up to maturity T
\end{itemize}
